\section{The nuclear fuel cycle}
The nuclear fuel cyclye encompasses the activities and processes 
for the use of fissile materials in fission nuclear reactors \cite{tsoulfanidis_nuclear_2013}. 
This can include the use of uranuim, thorium, or even plutonium as the 
fuel for the reactor, with uranium being the most widely used fuel for 
power reactors \hl{FIND CITATION}. This work focuses on a uranium based 
fuel cycle. The nuclear fuel cycle begins with the mining of uranium ores from the earth 
and ends with either the final disposal of the radioactive waste produced 
\cite{tsoulfanidis_nuclear_2013}. There are mutliple variations of 
the uranium-based fuel cycle that differ based on the type and design of 
the reactors deployed. Two primary variations on the nuclear fuel cycle 
include the once-through and uranium and plutonium recycling fuel cycles 
\cite{tsoulfanidis_nuclear_2013}. Figure \ref{fig:fuel_cycle} shows 
an example of all possible steps in a nuclear fuel cycle. 

\begin{figure}
    \centering
    \includegraphics[scale=0.8]{nuclear_fuel_cycle.jpg}
    \caption{A complete view of all options available for the nuclear 
    fuel cycle. Graphic from \protect\cite{noauthor_stages_2020}.}
    \label{fig:fuel_cycle}
\end{figure}

\subsection{Once-through fuel cycle}
A once-through fuel cycle is characterized by fuel only going through one 
pass through a reactor before being disposed of. In Figure \ref{fig:fuel_cycle},
this fuel cycle option is shown through the gray line from the storage 
step to the disposal step. None of thematerial flows from the Reprocessing 
Facility in Figure \ref{fig:fuel_cycle} are included in a once-through 
fuel cycle. This is the fuel cycle option currently employed in the 
United States. 
%\begin{figure}
    \centering
    \begin{tikzpicture}[node distance=1.5cm]
        \node (ore) [facility] {Uranium ore};
        \node (enrichment) [facility, right of=ore, xshift=1cm]{Enrich};
        \node (llwsink) [facility, below of=enrichment]{Depleted uranium};
        \node (fabrication) [facility, right of=enrichment, xshift=1cm]{Fuel};
        \node (reactor) [facility, above of=fabrication, xshift=2cm]{Power Plant};
        \node (spentfuel) [facility, below of=reactor, xshift=2cm]{Interim Storage};
        \node (sinkhlw) [facility, below of=spentfuel, xshift=2cm]{Geologic Repository};

        \draw [arrow] (ore) --  (enrichment); 
        \draw [arrow] (enrichment) -- (fabrication);
        \draw [arrow] (enrichment) -- (llwsink);
        \draw [arrow] (fabrication) |- (reactor);
        \draw [arrow] (reactor) -| (spentfuel);
        \draw [arrow] (spentfuel) -| (sinkhlw);

        \end{tikzpicture}
    \caption{Material flow for a once-through nuclear fuel cycle. All material is 
    disposed of after use in a reactor. Adapted from \protect\cite{wigeland_identification_2011}.}
    \label{fig:once_through_background}
\end{figure}

\subsection{Recycling fuel cycle}
A recycling fuel cycle is characterized by fuel going through additional 
passes through a reactor after a separtations step to remove fission 
products. There are a variety of methods to employ recycling in a 
fuel cycle. One distinction between recycling options is the number of 
times fuel is reprocessed. If fuel is reprocessed and burned in a reactor
for a finite number of times before being disposed of, then this recycling 
scheme is often referred to as a limited recycle. However, if the fuel is 
repeatedly reprocessed and burned in a reactor, then this is 
often referred to as a continuous reprocessing scheme. 
%\begin{figure} 
    \centering
    \begin{tikzpicture}[node distance=1.5cm]
        \node (ore) [facility] {Uranium ore};
        \node (enrichment) [facility, right of=ore, xshift=1cm]{Enrich};
        \node (llwsink) [facility, below of=enrichment]{Depleted uranium};
        \node (fabrication) [facility, right of=enrichment, xshift=1cm]{Fuel};
        \node (reactor) [facility, above of=fabrication, xshift=2cm]{Power Plant};
        \node (spentfuel) [facility, below of=reactor, xshift=2cm]{Interim Storage};
        \node (sinkhlw) [facility, below of=spentfuel, xshift=4cm]{Geologic Repository};
        \node (separation) [facility, below of=reactor, yshift=-1.5cm]{Separation};

        \draw [arrow] (ore) --  (enrichment); 
        \draw [arrow] (enrichment) -- (fabrication);
        \draw [arrow] (enrichment) -- (llwsink);
        \draw [arrow] (fabrication) |- (reactor);
        \draw [arrow] (reactor) -| (spentfuel);
        \draw [arrow] (spentfuel) -| node[anchor=west]{spent fuel}(sinkhlw);
        \draw [arrow] (spentfuel) |- (separation);
        \draw [arrow] (separation) -| (fabrication);
        \draw [arrow] (9,-1.0)-- node[anchor=north, text width=1.5cm]{fission products}(sinkhlw);
        
        \end{tikzpicture}
    \caption{Material flow for a once-through nuclear fuel cycle. All material is 
    disposed of after use in a reactor. Adapted from \protect\cite{wigeland_identification_2011}.}
    \label{fig:recycle_background}
\end{figure}

\subsection{Enrichment}
In the nuclear fuel cycle, enrichment increases the relative abundance of 
uranium-235 compared to uranium-238 in the fuel. The capacity of an 
enrichment facility is based on a number of quantities \cite{tsoulfanidis_nuclear_2013}:
\begin{subequations}
    \begin{equation}
        F = \text{mass (kg) of feed material per unit time}
    \end{equation}
    \begin{equation}
        P = \text{mass (kg) of product material per unit time}
    \end{equation}
    \begin{equation}
        T = \text{mass (kg) of tails produced per unit time}
    \end{equation}
    \begin{equation}
        x_f = \text{weight fraction of $^{235}$U in the feed stream}
    \end{equation}
    \begin{equation}
        x_p = \text{weight fraction of $^{235}$U in the product stream}
    \end{equation}
    \begin{equation}
        x_t = \text{weight fraction of $^{235}$U in the tails stream}
    \end{equation}
    \begin{equation}
        SWU = \text{separtive work units (SWU); the physical work required to separate the isotopes}
    \end{equation}
\end{subequations}

Each of these quantites can be related using the following equations:
\begin{subequations}
    \begin{equation}
        F = P + T
    \end{equation}
    \begin{equation}
        x_fF = x_pP + x_tT
        \label{eq:enrichment_2}
    \end{equation}
    \begin{equation}
        SWU = [P*V(x_p) +T*V(x_t) - F*V(x_f)]*t
    \end{equation}
    \text{in which:}
    \begin{equation}
        V(x_i) = (2x_i - 1)*\ln\left(\frac{x_i}{1-x_i}\right)
    \end{equation}
    \begin{equation}
        t = \text{a given time period}
    \end{equation}
    \label{eq:enrichment}
\end{subequations}

The equations in Eq. \ref{eq:enrichment} show that as the desired enrichment level 
of the product feed increases, the feed mass must also increase, assuming everything 
else is held constant. The equations also show that as the desired product 
enrichemnt level increases, the amount of \gls{SWU} capacity needed also increases. 
However, if the product mass decreases, then the required \gls{SWU} capacity 
also decreases. Therefore, changes to the product enrichment level and mass 
throughput have opposite effects on the amount of feed material and required 
\gls{SWU} capacity. Changes to both quantities that are in opposite directions 
(i.e. an increase in one and a decrease in the other) will not lead to a clear 
expectation to the change in feed material of \gls{SWU} capacity. The effect 
on these variables will depend on the magnitude of the two changes.

\section{Fuel cycle simulators}
\gls{NFC} simulators are computational tools to model the flow of materials
and the commissioning and decommissioning of facilities for a given fuel 
cycle. \gls{NFC} simulators can be used to evaluate the transition between fuel cycle 
options (e.g. from a once-through fuel cycle to one with recycling), the 
performance of one fuel cycle with a growth in the demand of nuclear power, 
and the effect of perturbations on a given fuel cycle \cite{piet_dynamic_2011}. 

Stuff about ideal simulators \cite{huff_next_2010,brown_identification_2016}.

A variety of fuel cycle simulators have been developed with different modeling 
capabilities. Some of the fuel cycle simulators developed including VISION 
\cite{yacout_visionverifiable_2006}, 
\gls{DYMOND} \hl{better citation}, and ORION \cite{gregg_analysis_2012}. 
VISION was developed primarily to meet the program objectives of the 
\gls{AFCI} based on the functionalities of \gls{DYMOND}-US \cite{yacout_visionverifiable_2006}.
VISION dynamically models the mass flow of material through the 
entire fuel cycle and calculates metrics based on the \gls{AFCI} program 
objectives. New functionalities added to VISION, compared with \gls{DYMOND}-US, 
include economic costing and isotope decay. VISION is divided into multiple 
modules to model the mass flow, calculate metrics, perform control, or 
integration. The software was developed using Powersim software \cite{powersim_powersim_nodate},
which provides a modular structure for mass flow to occur between modules. 

\gls{DYMOND} is a fuel cycle simulator developed by \gls{ANL} that combines 
multiple modeling paradigms, including system dynamics, discrete events, and
agent-based modeling \cite{feng_sensitivity_2020}.

performs depletion of meterial inputs to a rector facility using 
ORIGEN \cite{wieselquist_scale_2020}.

ORION 

\subsection{\Cyclus}
\Cyclus is a dynamic, open source agent-based fuel cycle simulator. Built 
in C++, \Cyclus uses only open source and freely available libraries to 
provide full access to all users and developers. The 
\Cyclus architecture treats materials and facilities discretely and allows 
for variable fidelity levels \cite{huff_fundamental_2016}. These attributes
of \Cyclus allow for the software to easily model any fuel cycle scenario.

\Cyclus uses the notion of an \textit{agent} to to represent different 
components in the simulated fuel cycle. Agents are 
defined using the \Cyclus \gls{API}, which 
defines and develops agents. Generic \glspl{API} are used to allow users 
to develop their own suite of agent libraries and use them within \Cyclus. 
The \glspl{API} anticipate the structure on information about a given library 
that is required by the core \Cyclus kernel. This allows them to facilitate 
information sharing between the plug-in library and the \Cyclus framework. 
In addition to the flexibility in libraries and agents, this framework 
also allows for flexibility in licensing and distribution of the user 
defined libraries. Libraries are loaded without changes to the \Cyclus 
kernel and without unwanted transfer of sensitive information. 

The agent-based modeling paradigm employed by \Cyclus allows agent level 
modeling, as opposed to system level modeling. This paradigm allows different 
fuel cycle facilities, such as a reactor and a fuel fabrication plant, to 
be defined independently but still interact with each other in the 
simulation. There are three main groups of agents within the \Cyclus 
architecture: facilities, institutions, and regions. Facilities are 
the individual units in the fuel cycle that implement technology, 
such as a fuel fabrication facility or a uranium mine. Institutions 
manage the facilities, similar to a company. Regions provide geographic 
and political context for the institutions and regions, and can be thought 
of as similar to individual nations. The types of agents are related using 
a parent-child hierarchy: regions are
parents of institutions, and institutions are parents of facilities. This 
structure requires that institutions are responsible for the deployment 
and decommissioning of the facilities. 

The \Cycamore module provides a variety of libraries that can be 
dynamically loaded into \Cyclus for use in a simualation 
\cite{carlsen_cycamore_2014,huff_fundamental_2016}. These libraries 
include mutliple facilities, two institutions, and one region as of 
\Cycamore 1.3. The two institutions interact with the region (the 
cycamore::GrowthRegion) to deploy facilities as needed to meet a 
commodity demand specified by the cycamore::GrowthRegion. 
The cycamore::DeployInst manually defines the number of each facilitiy 
to be deployed at each time step. These facilities are then registered with 
the cycamore::GrowthRegion for how they contribute to the demand. The 
facilities are decommissioned at the end of their lifetime. The 
cycamore::ManagerInst deploys facilities according to the commodity 
demand that is not met by facilities deployed by the cycamore::DeployInst. 

In \Cyclus, agents trade materials through the \gls{DRE} 
\cite{gidden_agent-based_2015,huff_fundamental_2016}. The \gls{DRE} defines the 
supply-demand communication framework and treats facilities as black boxes
so the solution strategy is agnostic to the resource types being exchanged. 
There are three steps to the \gls{DRE}: information gathering, resource 
exchanges, and trade execution. During the information gathering 
phase an exchange graph is created based on information from suppliers 
and consumers with any constraints. The requested demand from an agent 
can be more than what is actually consumed, granted capacity constraints 
accompany the request. An agent can request multiple commodities, but 
a preference of each commodity must be provided. A solution is then found to the 
exchange graph, then final trades are constructed and executed 
\cite{gidden_methodology_2016}. A feasible solution is found to the 
exchange graph using an optimization-based 
approach with econimics-based proxies \cite{gidden_methodology_2016}. 
The \gls{DRE} can be used to apply physical, economic, and social 
contraints to a fuel cycle model. 


SImialr to many of the other fuel cycle simulators developed, material 
compositions in \Cyclus are 
defined using recipes. However a module has been developed to allow for 
depletion calculations of a single enrichment reactor core using ORIGEN 
\cite{skutnik_cyborg:_2016}. This module increases \Cyclus's ability 
to use realistic material definitions for spent fuel. This is very 
important for an accurate understanding of reprocessing and repository 
needs. 

\subsection{Verification and use of fuel cycle simulators}
One key piece of work using fuel cycle simulators was a multi-lab effort 
to verify their results \cite{feng_standardized_2016}.

Verification of \Cyclus has been performed \cite{bae_standardized_2019}, 
based on a transition scenario from an open fuel cycle to an advanced
fuel cycle with reprocessing \cite{feng_standardized_2016}. Due to 
some inherent differences in \Cyclus and the other fuel cycle simulators 
used by \cite{feng_standardized_2016} there were some differences in 
the scenario parameters, such as \gls{LWR} batch size or 
cycle length. These differences are implemented to conserve the reactor 
core size, with a negligible difference in the \gls{SFR} core size. \Cyclus 
was shown to have strong agreement with the verification of other 
fuel cycle simulators \cite{bae_standardized_2019} with respect to 
multiple metrics such as reactor startup and decommission schedule, fresh 
fuel loading, and annual reprocessing throughputs. One major difference
between \Cyclus and other fuel cycle simulators identified by this work 
is that \Cyclus fully resolves discrete batches for fuel discharge. This 
implementation allows \Cyclus to better match realistic fuel handeling, but 
does not prevent \Cyclus from producing comparable results to other 
fuel cycle simulators. 

The use of \gls{NFC} simulators has often been to evaluate how fuel cycle options 
can meet specific goals. For example, \gls{DYMOND}, \gls{DANESS}, and 
\hl{NFCSim} were used to evaluate how different fuel cycle options 
could reduce \gls{SNF} volume in a repository \cite{yacout_dynamic_2000}. 

Use of \gls{NFC} simulators over the 
last decade has primarily been guided by the Nuclear Fuel Cycle \gls{ES} from the \gls{DOE-NE} 
\cite{wigeland_nuclear_2014}. This project categorized 
multiple fuel cycle options into 40 different \glspl{EG} based on characteristics 
such as the number of times spent fuel is recycled, the type of fuel burned, 
and the type of reactor (e.g. thermal critical reactor, fast critical reactor). 
The \gls{EG} were compared on their performance at equilibrium on nine evaluation 
criterion, including nuclear waste management and resource utilization. The 
results of this project showed that the most benefit comes from fuel cycles 
that use continuous recycle of U/Pu with new or natural uranium in fast critical 
reactors (\gls{EG} 23), continuous recycle of U/TRU with new or natural 
uranium in fast critical reactors (\gls{EG} 24), continuous recycle of U/Pu 
with new or natural uranium in both fast and thermal critical reactors 
(\gls{EG} 29) or contiuous recycle of U/TRU with new or natural uranium in 
both fast and critical reactors \cite{wigeland_nuclear_2014}. 

Based on the results of the \gls{ES}, many efforts have been made to model 
the transition to these promising fuel cycle options using fuel cycle 
simulators. ORION was used to model the transition from a once-through 
fuel cycle option using enriched uranium in thermal critical reactors (\gls{EG}
01) to \gls{EG} 29 \cite{sunny_transition_2015}. This work in ORION 
deomonstrated the importance of new \glspl{LWR} built or license extensions of 
exisiting reactors for this transition. 

\Cyclus has been used to model the \gls{EG} 23 transition scenario assuming 
a 1\% growth in power demand 
\cite{djokic_application_2015}. The results from \Cyclus were compared 
to the results from modeling the same transition scenario with \gls{DYMOND}, 
a fuel cycle simulator from \gls{ANL}. The power generated, reactor entry 
time, and reactor exit time for the \Cyclus simulation qualitatively 
matches well with the results of the \gls{DYMOND} simulation, but the power 
demand growth could match better to the target 1\% by slightly modifying 
the \gls{SFR} deployment. 

\section{Sensitivity analysis and optimization}
Sensitivity analysis is a broad termed used to encompass sensitivity analysis, 
uncertainty propagation, and parametric studies. 

Sensitivity analysis for fuel cycle simulations has been done using a variety of 
tools. Coupeling of \Cyclus with the Metaheuristic Opimization Tool (MOT) from 
\gls{ORNL} demonstrated \cite{feng_sensitivity_2020}. 

One tool for performing sensitivity analysis and optimization is 
Dakota \cite{adams_dakota_2019}

\subsection{Dakota}
Dakota was developed by \gls{SNL} as a systematic way for scientists and 
engineers to "obtain improved or optimal designs or understand sensitivity
or uncertainty using simulation-based models" \cite{adams_dakota_2019}.
Dakota can easily be coupled to a variety of codes in a ``black-box'' manner
to remove supplying Dakota with any of the other program's source code. 

Dakota has been coupled to fuel cycle simulators such as \gls{DYMOND} and \Cyclus 
\cite{chee_sensitivity_2019}. This coupeling was used to perform one-at-a-time, 
synergistic, and global sensitivity analysis on the transition from \gls{EG} 01
to \gls{EG} 30. 

Dakota has been used with \gls{DYMOND} to perform sensitivity analysis on a transition
scenario of "the US nuclear fleet to a closed fuel cycle with small modular 
\glspl{LWR} and fast reactors fueled by reprocessed used nuclear fuel" 
\cite{richards_application_2021}. This analysis focused on failure determination
and global analysis through the Sobol' indicies, and deomnstrated how the 
combination of these tools can ``identify important design and policy parameters
for transition scenarios and how those parameters interact not only with a single 
response, but also synergistically between all parameters and metrics of 
interest'' \cite{richards_application_2021}. 

\subsection{Optimization}
Optimization refers to the search for a solution of a given problem to 
meet specific objectives. When applied to fuel cycles, optimization 
is the determination of which fuel cycle option or facility 
deployment schedule will meet the desired goals. Common factors 
considered with optimization of a nuclear fuel cycle are natural 
resource utilization, economics, environmental impacts, and proliferation 
concerns \cite{passerini_systematic_2014,wigeland_nuclear_2014}. Optimizing 
a fuel cycle can be solved as a sinlge-objective or multi-objective 
problem. 

A variety of algorithms can be used to optimize a given 
fuel cycle. Linear programming optimization tools have been used for 
optimization, but have intrinsic issues for solving non-linear 
problems, such as fuel cycles, as demonstrated by \cite{passerini_systematic_2014}.
Linear programming optimization tools will often find a local minimum, 
but not the global minimum. 
Another method explored to optimize fuel cycle options is genetic 
algorithms \cite{passerini_systematic_2014}. Genetic algorithms 
have been shown to find a global minimum when multiple objectives 
are applied \cite{passerini_systematic_2014}. 


\section{HALEU}
\gls{HALEU} is defined as uranium that is enriched between 5-20\% by mass in 
uranium-235, and is considered a subset of \gls{LEU}. Fuel for \glspl{LWR} 
is enriched to no more than 5\% uranium-235, per \gls{NRC} regulations 
\hl{citation}. Table \ref{tab:enrichemnt} provides definitions of some of the 
main classifications of uranium enrichemnt. 

\begin{table}
    \centering
    \caption{Categories of uranium enrichment by weghit fraction of 
    uranium-235.}
    \label{tab:enrichemnt}
    \begin{tabular}{l c c}
        \hline
        Category & Weight fraction (\%)\\\hline
        Depleted & <0.711 \\
        Natural & 0.711 \\
        \gls{LEU} & 0.711-20 \\
        \gls{HALEU} & 5-20 \\
        \gls{HEU} & $\ge$20 \\
        \hline
    \end{tabular}
\end{table}

\subsection{Production Methods}
There are two primary methods of producing \gls{HALEU}. The first 
option is to enrich uranium to the required level. There is only one 
facility in the US currently 
licensed to enrich uranium, the Urenco LES facility in Eunice, 
New Mexico \cite{noauthor_establishing_2022}. This facility is only 
licensed to enrich uranium up to 5.5\% uranium-235 \cite{noauthor_establishing_2022},
which is less than many of the \gls{HALEU}-fueled reactor designs 
require. Another enrichment facility is currently under development in 
Piketon, Ohio which is licensed to enrich uranium up to 20\% 
uranium-235 \cite{noauthor_establishing_2022}. However, this is only 
designed to be a demonstration facility producing up to 600 kg of 
\gls{HALEU} \cite{us_nrc_centrus_2021} with the possibility of 
developing more facilities to produce up to 12 MTU/year of \gls{HALEU}
\cite{noauthor_establishing_2022}.

The other option is to downblend \gls{HEU} to the required \gls{HALEU}
level. The \gls{DOE} is considering muliple sources of \gls{HEU} that can 
be used to create \gls{HALEU} \cite{noauthor_establishing_2022}. The 
first is spent fuel from \gls{EBR}, which can produce a total of 10 MTU 
of \gls{HALEU} at 19.75\% enrichment \cite{noauthor_establishing_2022}. 
The next source is the stored \gls{HEU} solution at \gls{SRS} from 
spent research reactor fuel. This source is smaller than what is 
available from \gls{EBR}, and would only produce 2 MTU of \gls{HALEU} 
at 19.75\% enrichment \cite{noauthor_establishing_2022}. These sources of 
\gls{HALEU} would certainly be beneficial to fueling reactors, but they 
are a very limited supply of \gls{HALEU} that can not be counted on 
for long-term support of \gls{HALEU} reactors. 


\subsection{Fuel forms}
Advanced reactor designs that require \gls{HALEU} in a variety of fuel 
forms (Table \ref{tab:fuel_forms}). Each of the fuel forms required has 
distinct advantages and disadvantages. 

\begin{table}
    \centering
    \caption{Fuel form required by select advanced reactors that will 
    require \gls{HALEU}. Reactor designs listed here are all of US origin, 
    and information was obtained from \cite{hussain_advances_2018} unless 
    otherwise specified.}
    \label{tab:fuel_forms}
    \begin{tabular}{c c c}
        \hline
        Reactor Name & Company & Fuel form \\\hline 
        EM$^2$ & General Atomics & UC pellet \\
        eVinci & Westinghouse & UO$_2$ or UN \\
        Mk 1 PB-FHR & UC Berkeley & \gls{TRISO} pebbles\\
        \gls{MMR} \cite{mitchell_usnc_2020} & \gls{USNC} & UCO \gls{FCM} compact\\
        Natrium & Terrapower & Metallic \\
        SC-HTGR & Framatome & UCO \gls{TRISO} compacts \\
        Superstar  & \gls{ANL} & U/Pu metallic fuel \\
        Westinghouse Lead Fast Reactor  & Westinghouse & Oxide \\
        Xe-100 \cite{harlan_x-energy_2018} & X-energy & \gls{TRISO} pebbles \\
        \hline        
        
    \end{tabular}
\end{table}

\gls{TRISO} fuel particles contain 
a small fuel kernel surrounded by four layers of material: a porous carbon 
buffer layer, and inner pyrolytic carbon layer, a silicon carbide layer, 
and an outer pyrolytic carbon layer \cite{demkowicz_coated_2019}. The fuel 
kernel in the center contains the uranium fuel, typically in an oxide 
or oycarbide form. The porous carbon buffer acts an attenuator for any 
recoil fission products and accomodate pressure changes for gas accumulation.
The pyrolytic carbon layers protect the silicon carbide layer from 
any chemical degredation, while the silicon carbide layer maintains the 
structure integrity of the particle from any stresses and is the primary 
barrier against fission product gas release. The actual dimensions 
of each layer depends on the exact manufacturing and design purpose, but 
the \gls{TRISO} kernels are typically between 350-600 $\mu$m in diameter 
\cite{demkowicz_coated_2019}. \gls{TRISO} particles are known to have 
favorable performance at high temperatures 
\cite{demkowicz_coated_2019}. 

Multiple \gls{TRISO} kernels are combined together into a a pebble or 
a cyclindrical compact for placement in a graphite block, based on the 
design of the reactor \cite{demkowicz_coated_2019}. Each form of 
\gls{TRISO} particles has been used previously; the \gls{HTTR} uses 
\gls{TRISO} particles in cylindrical compacts in a graphite block 
\cite{shiozawa_overview_2004} and the \hl{find pebble example}.
\gls{TRISO} particles can also be placed in a silicon-carbide matrix 
to form \gls{FCM} fuel. This fuel form performs better in high-burnup 
applications because of improved stability under neutron irradtion and 
larger thermal conductivity \cite{snead_fully_2011}. The use of 
\gls{FCM} fuel in \glspl{LWR} to burn transuranic material from waste 
has been investigated \cite{snead_fully_2011,venneri_fully_2011}. 

\subsection{Expected Demand}
A survey of select advanced reactor companies by \gls{NEI} reports an 
expected cumulative demand of \gls{HALEU} of almost 600 MT by 2030 
\cite{korsnick_need_2018}.

Demand for \gls{ARDP} reactors will begin in 2024, so that fuel is 
properly fabricated in time for reactors to come on-line \cite{noauthor_establishing_2022}.
Other companies are expecting to deploy working \gls{HALEU}-fueled reactors 
by the mid-2020's \cite{noauthor_establishing_2022}.

An estimated 20 MTU of \gls{HALEU} will be needed between 2024-2027 to 
meet the needs of already announced projects in the US with a site 
selected \cite{noauthor_establishing_2022}. An estimated 6 MTU/year of 
\gls{HALEU} will be needed in 2028 to support refueling of these projects 
\cite{noauthor_establishing_2022}. 

\subsection{HALEU in reactors}
A lot of work has been done to evaluate the performance of reactors 
using \gls{HALEU} fuel. Burns et al. investigated the reactor and fuel cycle 
performance of an \gls{LWR} fueled by an enrichment above 5\% \cite{burns_reactor_2020}.
Their work showed that increasing the fuel enrichment up to 7\% does not 
largely affect the fuel temperature or the moderator temperature coefficents,
but it does decrease the soluble boron coefficient and increase the maximum 
burnup ar the edges of the fuel pellets. The impacts on the fuel cycle include 
an increase in the amount of natural uranium required, a decrease in the 
amount of high-level waste disposed of per unit energy, and changes in the 
waste radioacitivity waste. 

Increasing the fuel enrichment of the NuScale \gls{SMR} has also been investigated 
\cite{carlson_implications_2022}. Increasing the fuel enricment to 8.34\% doubled the 
fuel discharge burnup and the cycle time. Fission product poisons, specifically 
sumarium-149 and xenon-135, increased on concentration when \gls{HALEU} fuel was used. 
This increase impacts the neutrons poisons that are needed during operation and the 
radioactivity of the fuel after discharge. Using \gls{HALEU} in the core also led 
to a reduction in the \gls{LCOE} for the reactor for certain combinations of 
enrichment and cycle length \cite{carlson_economic_2020,carlson_implications_2022}.

If \gls{HALEU} fuel is produced by downblending \gls{HEU}, is it very possible that 
the \gls{HALEU} will contain impurities that would not be present if natural uranium 
were enriched to produce the \gls{HALEU} \cite{noauthor_establishing_2022}. Bounding 
studies on the uranium isotopic 
composition of \gls{HALEU} produced from spent fuel from \gls{EBR} shows the presence 
of uranium-232, 233, 234, and 236 \cite{vaden_isotopic_2018}. \gls{HALEU} created from 
downblened \gls{HEU} at the Y-12 Nuclear Security Complex also has uranium-232, 234, and 
236 impurities \cite{nelson_foreign_2010}. Uranium-233 is a fissile 
isotope and is expected to affect the neutron population inside a reactor. Uranium-232, 
234, and 236 are parasitic neutron absorbers, so their presence is also expected to affect 
the neutron population. The neutronics performance of using \gls{HALEU} containing
uranium-235 and 238 was compared with the performance using \gls{HALEU} from Y-12 with the 
known impurities \cite{celikten_effects_2021}. Using 
the \gls{HALEU} with the \gls{HEU} impurities showed a 0.43\% reactivity drop, leading to 
an 8\% reduction in the cycle length. 
